Mapping high-resolution surface water in dynamic deltas is critical for water resource management. Most empirical machine-learning approaches map spectral reflectances directly to water masks, bypassing physical catchment state representations. We evaluate geomorphological and temporal controls on surface water in the Indus River Delta (Pakistan) from 2018 to 2024. Using 83,000 monthly GEE records, we analyze Random Forests, LSTMs, and Transformers under a leakage-free input configuration. Sequence models exhibited collapse under the evaluated coordinate-withheld setting during dry periods ($F_1\text{-score} = 0.000$) due to scale mismatches between coarse climate forcings and sub-pixel channels. During extreme floods, spatial Random Forests outperform sequence models ($F_1\text{-score} = 0.845$ vs. $0.745$) because static features constrain water boundaries. Ablation experiments show that a Terrain-only model achieves an $F_1\text{-score}$ of 0.842 (95\% CI: [0.833, 0.851]), which shows no statistically significant difference from the Combined model ($F_1\text{-score} = 0.845$; McNemar's $p = 0.2584$, representing limited statistical power to detect small differences rather than proof of equivalence). These results indicate that high-resolution geomorphological constraints explain the spatial organization of inundation boundaries under the evaluated feature set, whereas coarse meteorological forcing primarily regulates temporal activation. Furthermore, lag correlation audits identify a 3-month upstream precipitation routing delay in coastal wetlands and evaluate a negative correlation in croplands ($r = -0.299$) that remains ambiguous due to confounding crop phenology and irrigation signals. These findings motivate the future development of physically constrained hydrological state-space models.
